\documentclass[12pt,a4paper]{report}
\usepackage[T1]{fontenc}
\usepackage[utf8]{inputenc}
\usepackage{times}
\usepackage[top=1in, bottom=1in, left=1.2in, right=1in]{geometry}
\usepackage{setspace}
\onehalfspacing
\usepackage{graphicx}
\usepackage{amsmath,amssymb}
\usepackage{hyperref}
\usepackage{titlesec}
\usepackage{float}
\usepackage{longtable}
\usepackage{booktabs}
\usepackage{listings}
\usepackage{xcolor}

\definecolor{codegreen}{rgb}{0,0.6,0}
\definecolor{codegray}{rgb}{0.5,0.5,0.5}
\definecolor{codepurple}{rgb}{0.58,0,0.82}
\definecolor{backcolour}{rgb}{0.95,0.95,0.92}

\lstdefinestyle{mystyle}{
    backgroundcolor=\color{backcolour},   
    commentstyle=\color{codegreen},
    keywordstyle=\color{magenta},
    numberstyle=\tiny\color{codegray},
    stringstyle=\color{codepurple},
    basicstyle=\ttfamily\footnotesize,
    breakatwhitespace=false,         
    breaklines=true,                 
    captionpos=b,                    
    keepspaces=true,                 
    numbers=left,                    
    numbersep=5pt,                  
    showspaces=false,                
    showstringspaces=false,
    showtabs=false,                  
    tabsize=2
}
\lstset{style=mystyle}

\title{\textbf{Survival Guide for MS Economics Dissertation Defense}\\
\large A Beginner-Friendly, Step-by-Step Breakdown for Engineering Students}
\author{Your Personal AI Tutor}
\date{\today}

\begin{document}

\maketitle
\tableofcontents
\newpage

\chapter{Introduction \& The Big Picture}
\section{What is this Dissertation About?}
Your dissertation is titled: \textbf{Do Global Oil Price Shocks Raise India's Inflation More Than They Lower It? Short-Run Pass-Through to CPI Inflation in India, 2004--2024}. 

As an engineering student transitioning to MS-level economics, all the jargon can be intimidating, but the core idea of your thesis is actually very simple and intuitive. Imagine you run a logistics company. If the price of diesel goes up, you immediately raise your shipping fees to cover costs. But if the price of diesel goes down a few months later, do you immediately lower your shipping fees? Probably not right away, or by as much. You like keeping the extra profit margin! 

This phenomenon—where prices go up fast like a \textbf{rocket} but come down slowly like a \textbf{feather}—is called \textbf{asymmetric pass-through}. Your dissertation simply investigates if this happens at a national level in India. India imports about 85\% of its crude oil. When global oil prices jump, does Indian consumer inflation spike more severely than it drops when oil prices fall? 

\section{If the Professor Asks: "Why this topic?"}
\textbf{Professor:} "Why did you choose this exact topic?" \\
\textbf{You:} "I wanted to study how global shocks transmit to domestic prices. Since India imports 85\% of its oil, crude oil is the perfect transmission channel. I chose to focus on \textit{asymmetry} because classical economic models assume prices adjust symmetrically. However, real-world rigidities (like sticky prices, government interventions, or firm profit-seeking) suggest prices are 'sticky downward'. Testing this asymmetry using the NARDL model gives us a much more realistic picture of inflation dynamics in India over the past 20 years."

\section{Key Economic Terms You Must Know}
\begin{itemize}
    \item \textbf{CPI (Consumer Price Index):} A measure of the average change over time in the prices paid by urban and rural consumers for a market basket of consumer goods and services. \textit{In simple terms:} It's the official measure of inflation. If CPI goes up, the cost of living went up.
    \item \textbf{Brent Crude:} A major trading classification of sweet light crude oil. It serves as a major benchmark price for purchases of oil worldwide. \textit{In simple terms:} It's the standard global price of raw oil. 
    \item \textbf{IIP (Index of Industrial Production):} An index that tracks manufacturing activity in the economy. You use this as a 'control variable' to account for general economic demand.
    \item \textbf{Exchange Rate (EXR):} The value of one currency for the purpose of conversion to another. You use INR/USD. Since oil is priced in USD globally, but Indians buy fuel in Rupees (INR), the exchange rate heavily affects how much an oil shock actually hurts India.
    \item \textbf{Pass-Through:} How much of a price change at the source (global oil) is passed down to the final consumer (CPI). 
\end{itemize}

\newpage
\chapter{The Data \& Python Pre-Processing}
\section{What is Time Series Data?}
A time series is simply data collected at regular intervals over time. Your criteria required 19-20 years of data. You used monthly data from \textbf{April 2004 to December 2024} (249 months, which is 20.75 years). This perfectly satisfies the professor's requirement.

\section{The Python Script: \texttt{chain\_link\_iip.py}}
Why did you write a Python script? Data collected over 20 years rarely stays perfectly consistent. The Indian government measures IIP, but they update their "base year" (the reference year set to 100) every decade or so. 
Your data spans two different base years for IIP: \textbf{Base 2004-05} and \textbf{Base 2011-12}. You cannot just paste the two columns together; a value of 150 in the 2004 base does not equal 150 in the 2011 base.

\textbf{What the Python script does:}
It performs "Chain-Linking" using the \textbf{Overlap-Ratio Method}. Since there is a period (April 2012 to Jan 2017) where the RBI reported the IIP in \textit{both} bases, the script calculates a "splice factor" (the average ratio between the two bases during the overlap). It then multiplies the older 2004 base data by this factor to smoothly convert it into the newer 2011 base currency equivalent without creating any artificial jumps.

\textbf{Professor:} "How did you merge the IIP data?"\\
\textbf{You:} "I used the overlap-ratio chain-linking method. I took the period where both bases (2004-05 and 2011-12) were available, calculated the average ratio between them (the splice factor, which came to approximately 0.627), and scaled the old base to match the new base. I plotted the final series to visually verify there was no structural break at the splice point."

\newpage
\chapter{Time Series Econometrics: The Fundamentals}
As an engineer, you know calculus and statistics. Econometrics is just applying statistics to economic data to find causal or predictive relationships. But Time Series has special rules because data points are ordered in time.

\section{Stationarity and the Unit Root Test}
\textbf{The Concept:} If you run a standard linear regression (like $y = mx + c$) on two variables that are just growing over time (like human population and your age), the math will wrongly tell you they are highly correlated just because they both go up over time. This is called a \textbf{spurious regression}.

To prevent this, variables need to be \textbf{Stationary} (mean and variance are constant over time). If a variable is not stationary, we say it has a \textbf{Unit Root} (it's integrated of order 1, or $I(1)$). To make it stationary, we usually take the "first difference" (i.e., subtract yesterday's value from today's value, $\Delta y_t = y_t - y_{t-1}$).

\textbf{The Test:} You used the \textbf{ADF (Augmented Dickey-Fuller) Test}.
\begin{itemize}
    \item \textbf{Null Hypothesis ($H_0$):} The data has a unit root (is non-stationary).
    \item \textbf{Alternative ($H_1$):} The data is stationary.
\end{itemize}
If the p-value is $< 0.05$, we reject $H_0$ and conclude the data is stationary. 

\textbf{In your thesis:} You took the natural logarithm of CPI, Oil, IIP, etc. The ADF test showed that the log-levels were non-stationary, but their first differences (the month-over-month percentage changes, like $\Delta \ln CPI$) were strongly stationary ($p < 0.01$). This allowed you to use these variables safely in regressions.

\section{Lags}
A "Lag" is simply a past value. Because it takes time for a global oil price increase to reach Indian retail prices, we must include "lags" of the oil price in our model. Over your model, you use up to $L=4$ (4 months back) for CPI and $L=3$ (3 months back) for oil.

\newpage
\chapter{The Main Model: NARDL (Nonlinear ARDL)}

\section{What is ARDL?}
ARDL stands for \textbf{Autoregressive Distributed Lag}.
\begin{itemize}
    \item \textbf{Autoregressive (AR):} Uses past values of the dependent variable (CPI) to predict its current value. Inflation often has inertia.
    \item \textbf{Distributed Lag (DL):} Uses present and past values of the independent variables (Oil, IIP, Exchange rate).
\end{itemize}

\section{Why 'Nonlinear' / Asymmetric? (NARDL)}
A standard ARDL lumps all oil price changes into one variable. But we want to test rockets and feathers! So, we create two new variables:
\begin{itemize}
    \item $\Delta Oil^+$: Contains only the magnitude of oil price \textit{increases}. If the price drops, this is set to 0.
    \item $\Delta Oil^-$: Contains only the magnitude of oil price \textit{decreases}. If the price goes up, this is zero.
\end{itemize}
By putting both into the regression separately, the model gives us two separate coefficients: $CPT^+$ (Cumulative Pass-Through for positive shocks) and $CPT^-$ (Cumulative Pass-Through for negative shocks).

\section{Dummy Variables}
You also included "Dummies" (variables that are just 1 or 0) to control for major policy regime changes and extreme outlier events:
\begin{itemize}
    \item \textbf{D\_petrol:} Set to 1 after June 2010 when gasoline prices were deregulated.
    \item \textbf{D\_diesel:} Set to 1 after Oct 2014 when diesel was deregulated.
    \item \textbf{D\_covid:} Set to 1 for April 2020 to absorb the extreme lockdown shock.
    \item \textbf{Seasonal Dummies (M1-M11):} To remove predictable monthly seasonality.
\end{itemize}

\textbf{Professor:} "Why use NARDL instead of VAR or VECM?"\\
\textbf{You:} "NARDL allows us to cleanly split our shock variable into positive and negative partial sums, directly estimating the asymmetry in short-run dynamics and long-run cointegration. Furthermore, ARDL models are extremely robust even when dealing with a mix of I(0) and I(1) variables, which avoids the strict pre-testing requirements of the Johansen cointegration approach used in VECM."

\newpage
\chapter{Understanding the R Code (\texttt{analysis.R})}
If the professor asks you to "walk through the code" or "run it on the software", do not panic. The R script is just a step-by-step calculator of everything we discussed. 

\section{Step 1: Data Loading}
The script uses standard R tools like \texttt{tidyverse} and \texttt{readxl} to load the CSVs (from FRED) and your newly chained IIP excel file. It merges them by date and crops them precisely to your window: \texttt{2004-04-01} to \texttt{2024-12-01}.

\section{Step 2: Variable Construction}
\begin{lstlisting}[language=R]
# Convert global brent USD to INR
df$oil_inr <- df$brent_usd * df$exr

# Take natural logs
df$ln_cpi <- log(df$cpi)
df$ln_oil <- log(df$oil_inr)

# Take differences (which approximations percentage change)
df$dlnCPI <- c(NA, 100 * diff(df$ln_cpi))
df$dlnOil <- c(NA, 100 * diff(df$ln_oil))

# The "Rockets and Feathers" secret sauce!
df$dlnOil_pos <- ifelse(!is.na(df$dlnOil), pmax(df$dlnOil, 0), NA)
df$dlnOil_neg <- ifelse(!is.na(df$dlnOil), pmin(df$dlnOil, 0), NA)
\end{lstlisting}

\section{Step 4 \& 6: AIC Lag Selection & The Main Regression}
How many past months (lags) should we look at? We don't guess. The script loops over 1 to 4 lags and uses \textbf{AIC (Akaike Information Criterion)}. AIC calculates a score for how well the model fits penalized by how complex it is. The lower the AIC, the better. The script finds the optimal number of AR lags is $p=1$.

Then it runs the OLS linear regression: \texttt{lm(formula, data)}.
\begin{lstlisting}[language=R]
asym_model <- lm(dlnCPI ~ dlnCPI_L1 + 
                 dlnOil_pos_L0 + dlnOil_pos_L1 + ... +
                 dlnOil_neg_L0 + dlnOil_neg_L1 + ... +
                 dlnIIP + D_petrol + D_diesel + D_covid + M1...M11)
\end{lstlisting}

\newpage
\chapter{Diagnostic Tests: Making Sure the Model is Valid}
The professor will definitely ask, "Did you check the assumptions of your model?" Linear regression requires error terms (residuals) to behave randomly. If they don't, your p-values are trash. 

The R script automatically performs 4 standard tests on your model:

\section{1. Breusch-Godfrey LM Test}
\textbf{What it tests:} Autocorrelation (are errors from this month correlated with errors from last month?).\\
\textbf{Interpretation:} Null hypothesis ($H_0$) is \textit{no autocorrelation}. Your model passes ($p > 0.05$).

\section{2. Breusch-Pagan Test}
\textbf{What it tests:} Heteroskedasticity (does the variance/spread of the errors change over time?).\\
\textbf{Interpretation:} Null hypothesis ($H_0$) is \textit{homoskedasticity} (constant variance). If $p < 0.05$, we have heteroskedasticity. \\
\textbf{YOUR SAVING GRACE:} Your script detects this and actively corrects it! You calculate the "Newey-West HAC (Heteroskedasticity and Autocorrelation Consistent) Standard Errors". \\
\textbf{If the Prof asks:} \textit{"I used Newey-West standard errors to ensure my t-tests and F-tests remain totally valid even if heteroskedasticity is present in the residuals."}

\section{3. Ramsey RESET Test}
\textbf{What it tests:} Model mis-specification (did we miss a non-linear term?).\\
\textbf{Interpretation:} Null hypothesis ($H_0$) is \textit{correct specification}. Passes for your model.

\section{4. CUSUM Test}
\textbf{What it tests:} Parameter stability over time.\\
\textbf{Interpretation:} It plots cumulative sums of recursive residuals. If the line stays within the red boundary bands, the model's coefficients are politically and economically stable over the 20 years.

\newpage
\chapter{Results \& How to Interpret Them}

\section{The Core Research Question}
We wanted to know if $CPT^+$ (Cumulative Pass Through of positive shocks) is greater than $CPT^-$ (negative shocks).

The script outputs these exact values:
\begin{itemize}
    \item \textbf{CPT+ ($\approx 0.021$):} A 10\% increase in oil prices raises India's monthly CPI by about 0.21 percentage points.
    \item \textbf{CPT- ($\approx 0.00$):} A 10\% decrease in oil prices basically does nothing to drop inflation.
\end{itemize}

\textbf{Conclusion based on point estimates:} YES! We have clear rockets and feathers asymmetry structurally. Positive shocks raise inflation, negative shocks do not lower it.

\section{The Wald Test (The Twist!)}
However, in statistics, it's not enough for 0.021 to just be numerically bigger than 0. We have to prove they are statistically, significantly different taking error margins into account.

We run a \textbf{Wald Test} where $H_0$: $CPT^+ = CPT^-$.
The test gives a p-value of $\approx 0.24$. Since $0.24 > 0.05$, we fail to reject the null hypothesis.

\textbf{Professor:} "Your Wald test fails to reject symmetry. Does that mean your entire thesis is useless?"\\
\textbf{You:} "Not at all, sir. The point estimates clearly exhibit an asymmetric 'rockets and feathers' magnitude—positive shocks strictly pass through (0.21 pp) while negative shocks die out (~0 pp). However, the statistical variance is too large to definitively claim asymmetry at the strict 5\% significance level on the \textit{aggregate, nationwide headline CPI}. Headline CPI includes food (which is ~46\% of the basket and highly volatile). When I ran the exact same robustness check in the appendix using only the 'Fuel and Light' CPI sub-index, the asymmetry was dramatically larger and strictly statistically significant (p = 0.03). This proves the asymmetry exists strongly at the energy level, though it gets noisily diluted when diluted into the massive headline CPI basket."

\newpage
\chapter{Robustness Checks \& Advanced Questions}

Robustness checks show the professor that you didn't just run one model and pray. You tortured the data to see if your conclusions broke. They didn't. 

\textbf{1. Sub-sample Pre/Post 2014}\\
India fully deregulated diesel prices in October 2014. Before this, government subsidies heavily masked oil shocks. Your code literally splits the data at Oct 2014 and runs the models separately. \textit{Result:} The positive pass-through is much stronger and faster in the post-2014 deregulated era, exactly as economic theory predicts!

\textbf{2. Separating Brent USD from the Exchange Rate}\\
Usually, we multiplied Brent by INR/USD. But what if an oil shock is purely because the Rupee crashed, not because global oil went up? Your code estimates a model keeping Brent (in USD) and EXR separate. \textit{Result:} The Brent-specific positive pass-through remains clearly positive.

\textbf{3. Removing the COVID-19 dummy}\\
COVID was an anomalous mega-shock. What if it is skewing everything? Model re-run without the dummy. \textit{Result:} Pass-through characteristics barely change.

\textbf{4. Winsorizing at 1\%}\\
Replaces extreme top 1\% and bottom 1\% outliers with the 99th and 1st percentile values. \textit{Result:} No major change, proving results aren't just driven by one crazy month in 2008 or 2020.


\vspace{2cm}
\begin{center}
\textbf{\Large YOU ARE READY!} \\
\vspace{0.5cm}
Take a deep breath. You are an engineer who conquered a Master's level Economics dissertation using time-series econometrics in R and Python. Speak slowly, refer to the coefficients logically, and you will ace it. Good luck!
\end{center}

\end{document}
